\chapter{Literature Review, Innovation and Gaps / Revisão de Literatura, Inovação e Lacunas}

\section*{English}
The literature is organized into three relatively mature but weakly integrated streams: conformal broadband antennas and arrays; functional radomes and frequency-selective surfaces; and passive detection, localization and emitter identification. Conformal direction-finding studies show that mutual coupling, supporting structure and radome distortion belong to the calibrated array manifold, not merely to a mechanical enclosure \citep{caratelli2011,liberal2011}.

Functional-radome research demonstrates passband transmission combined with absorption or diffusion outside selected bands \citep{costa2012,lv2021,omar2018,sheng2024}. This creates a design tension for open spectrum monitoring: a radome optimized for stealth or narrow transmission can reject signals that a surveillance receiver should observe.

Passive-localization research uses AOA, TDOA and FDOA, with accuracy governed by geometry, synchronization, propagation and computational load \citep{liu2019,pine2021}. RF fingerprinting and open-set anomaly detection can assist emitter characterization, but closed-set classifiers and limited training conditions cannot be treated as universal solutions \citep{soltanieh2020,jagannath2022,sun2022}.

The project gap is the end-to-end integration of a conformal multiband aperture, distributed passive reception, polarimetric and temporal calibration, hierarchical event processing and realistic mountain-site validation. The innovation is therefore systemic and must be demonstrated with reproducible measurements.

\section*{Português}
A literatura é organizada em três linhas relativamente maduras, porém pouco integradas: antenas e arrays conformais de banda larga; radomes funcionais e superfícies seletivas em frequência; e detecção passiva, localização e identificação de emissores. Estudos de radiogoniometria conformal mostram que acoplamento mútuo, estrutura de suporte e distorção do radome pertencem ao manifold calibrado do array, não sendo apenas questões mecânicas \citep{caratelli2011,liberal2011}.

A pesquisa de radomes funcionais demonstra transmissão em banda combinada com absorção ou difusão fora das bandas selecionadas \citep{costa2012,lv2021,omar2018,sheng2024}. Isso cria uma tensão de projeto no monitoramento aberto do espectro: um radome otimizado para baixa assinatura ou transmissão estreita pode rejeitar sinais que o receptor deveria observar.

A pesquisa de localização passiva utiliza AOA, TDOA e FDOA, com precisão governada por geometria, sincronização, propagação e carga computacional \citep{liu2019,pine2021}. Fingerprinting de RF e detecção de anomalias em conjunto aberto podem auxiliar a caracterização de emissores, mas classificadores fechados e condições limitadas de treinamento não devem ser tratados como soluções universais \citep{soltanieh2020,jagannath2022,sun2022}.

A lacuna do projeto é a integração ponta a ponta de uma abertura conformal multifaixa, recepção passiva distribuída, calibração polarimétrica e temporal, processamento hierárquico de eventos e validação realista em sítios montanhosos. A inovação é sistêmica e deve ser demonstrada por medições reproduzíveis.
